% Requires \usepackage{booktabs} and \usepackage{multirow} in the preamble.
\subsection{Screening Candidate Predictive Models by Category}
\label{sec:hull_model_screening}

GeoSMILE's SMILE-surrogate attribution machinery is model-agnostic: it perturbs
inputs and observes the resulting shift in a model's own prediction
distribution, so it can explain any regressor exposing a plain
\texttt{fit(X, y)} / \texttt{predict(X)} interface. XGBoost was adopted as the
predictive model this thesis explains, but that choice was not asserted
without comparison. Before committing to it, six further candidates, spanning
the three model categories realistically available for a tabular,
non-self-explainable crime-count regression, were fit on the same panel and
temporal split and screened on predictive fit and on how faithfully GeoSMILE
could explain each one. The categories were: (i) \emph{gradient-boosted
trees} -- LightGBM and CatBoost, alongside XGBoost itself; (ii)
\emph{kernel/instance-based regressors} -- Support Vector Regression (SVR,
RBF kernel) and Gaussian Process Regression; and (iii) \emph{feed-forward
neural networks} -- a scikit-learn MLP and a deeper Keras/TensorFlow network,
representing the two neural-network toolkits most commonly used for this
class of problem. None of the six alternatives were hyperparameter-tuned;
each was sized to a comparable capacity to XGBoost's published configuration,
since the purpose of this pass is screening, not optimisation.

Every candidate was trained twice, once per branch (Feature: 30 non-spatial
predictors; Geo: the two LSOA centroid coordinates), with the same seed (42),
the same temporal train/validation/test split, and the same SMILE budget as
the published design (Feature: 1{,}024 masks; Geo: 4{,}000 perturbations --
raised from the originally published 2{,}000 after a further XGBoost sweep
showed higher Geo fidelity at that budget). A guard was applied before
committing to the expensive attribution step: a branch was fit and explained
by GeoSMILE only if its held-out predictive $R^2$ exceeded 0.7; a branch at or
below that threshold was excluded from GeoSMILE fitting and every downstream
attribution metric, rather than explaining a model that cannot itself
predict adequately. This screen deliberately excludes any metric benchmarked
against the external expert-reference set (ATT accuracy/F1/AUROC): those
measure agreement with a human-curated list, not whether GeoSMILE can
faithfully and stably explain a given predictive model, which is the
question this comparison is asking. Table~\ref{tab:hull_model_screening}
therefore reports, per branch, only: held-out predictive $R^2$; the
Dehghani-style held-out surrogate fidelity $R^2$ (evaluation-split, not the
kernel-selection diagnostic or the consistency-run mean -- see
\S\ref{sec:hull_fidelity_definitions}); rank-Spearman consistency of the
attribution ranking over ten SMILE reseeds ($R=10$); and cardinal-neighbour
ATT stability (top-$k$ Jaccard of the attribution set before/after adding one
real neighbouring LSOA to the reference).

\begin{table}[htbp]
\centering
\caption{Predictive fit and GeoSMILE attribution quality by candidate model and category (seed 42, Feature $N{=}1{,}024$, Geo $N{=}4{,}000$). ``--'' marks a joint Feature+Geo diagnostic (stability) withheld because the model's Geo branch failed the $R^2>0.7$ guard.}
\label{tab:hull_model_screening}
\resizebox{\textwidth}{!}{%
\begin{tabular}{llrrrrrrrr}
\toprule
& & \multicolumn{4}{c}{Feature} & \multicolumn{4}{c}{Geo} \\
\cmidrule(lr){3-6}\cmidrule(lr){7-10}
Category & Model & $R^2$ & Fidelity & Consist.\ $\rho$ & Stability $J$ & $R^2$ & Fidelity & Consist.\ $\rho$ & Stability $J$ \\
\midrule
\multirow{3}{*}{Gradient boosting}
 & XGBoost  & 0.8957 & 0.9780 & 0.647 & 1.000 & 0.8752 & 0.9479 & 0.945 & 1.000 \\
 & LightGBM & 0.8968 & 0.9834 & 0.722 & 1.000 & 0.8798 & 0.9472 & 0.950 & 1.000 \\
 & CatBoost & 0.8959 & 0.9849 & 0.855 & 1.000 & 0.8739 & 0.9482 & 0.943 & 1.000 \\
\midrule
\multirow{2}{*}{Kernel / instance-based}
 & SVR & 0.8629 & 0.9366 & 0.823 & -- & 0.3161 & \multicolumn{3}{c}{excluded, $R^2 \le 0.7$} \\
 & GP  & 0.8993 & 0.9964 & 0.955 & 1.000 & 0.8777 & 0.9461 & 0.961 & 1.000 \\
\midrule
\multirow{2}{*}{Neural network}
 & MLP        & 0.8529 & \textbf{0.0905} & 0.428 & -- & 0.4444 & \multicolumn{3}{c}{excluded, $R^2 \le 0.7$} \\
 & DL (Keras) & 0.8574 & \textbf{0.1244} & 0.586 & 1.000 & 0.8235 & 0.9451 & 0.953 & 1.000 \\
\bottomrule
\end{tabular}%
}
\end{table}

\paragraph{Gradient-boosted trees.} XGBoost, LightGBM, and CatBoost are
essentially interchangeable on predictive fit (Feature $R^2 = 0.896$--$0.897$;
Geo $R^2 = 0.874$--$0.880$) and on GeoSMILE fidelity (Feature $\ge 0.978$, Geo
$\ge 0.947$ across all three), with perfect cardinal-neighbour stability
($J=1.000$) on both branches for all three. Consistency is where they
separate: CatBoost's Feature-branch ranking is the most stable under reseeding
($\rho=0.855$ at $R{=}10$), against $0.647$ for XGBoost and $0.722$ for
LightGBM, while all three cluster tightly on the Geo branch
($\rho=0.943$--$0.950$). This is the only category in which every tested
model cleared the $R^2 > 0.7$ guard on both branches.

\paragraph{Kernel and instance-based regressors.} Restricted to these four
metrics, Gaussian Process regression is not a weak alternative -- it is the
strongest single candidate in the table on Feature-branch quality: the best
predictive fit of any model screened ($R^2=0.899$) and the best Feature-branch
fidelity by a wide margin ($R^2=0.996$, next-best CatBoost at $0.985$), with
the best Geo-branch consistency of any model that clears the $R^2>0.7$ guard
on that branch ($\rho=0.961$; the MLP's nominal $0.988$ is not a valid
comparison, since its Geo branch fails the guard outright, see below).
GP's stability is perfect on both branches. Nothing in
Table~\ref{tab:hull_model_screening} disqualifies it on quality grounds. What
does not appear in the table is cost: across the full screening protocol
(main pass, four cardinal-neighbour stability directions, ten consistency
reseeds) GP alone took approximately four hours to explain, against
single-digit minutes for any boosted-tree candidate, because Gaussian
Process prediction cost scales with training-set size and GeoSMILE issues
thousands of individual predictions per branch. GP is therefore carried
forward as the strongest quality case for an alternative to boosting, with
computational cost as the open objection, rather than being ruled out on
these four metrics. SVR illustrates the guard directly: its Feature branch
is workable ($R^2=0.863$, fidelity $0.937$, $\rho=0.823$) but its Geo branch,
given only two raw coordinates and no spatial kernel structure, could not
predict crime counts above $R^2=0.316$ and was excluded from GeoSMILE
fitting entirely; because stability is a joint Feature+Geo diagnostic, it is
withheld for SVR's Feature branch too, even though that branch individually
passed the guard.

\paragraph{Neural networks.} Both the scikit-learn MLP and the deeper
Keras/TensorFlow network predict the Feature branch adequately in isolation
($R^2=0.853$ and $0.857$) but GeoSMILE's surrogate fidelity on that same
branch collapses for both: $R^2=0.090$ for the MLP and $0.124$ for the deep
network, both negative once the weighted-adjustment for held-out sample size
is applied ($-0.073$ and $-0.008$ respectively). Feature-branch consistency
is correspondingly the weakest in the whole screen for these two models
($\rho=0.428$ and $0.586$). The two toolkits diverge on the Geo branch,
however: the deep network predicts it adequately ($R^2=0.824$) and GeoSMILE's
fidelity ($0.945$) and consistency ($\rho=0.953$) there are comparable to, or
better than, the boosted trees, with perfect stability. The MLP's Geo branch,
by contrast, fails the same $R^2>0.7$ guard that excludes SVR's Geo branch
($R^2=0.444$) and is excluded from GeoSMILE fitting entirely -- because
stability is a joint Feature+Geo diagnostic, it is withheld for the MLP's
Feature branch too, even though that branch individually passed the guard,
mirroring SVR's treatment. Predictive fit alone would not have flagged either
network as unsuitable on the Feature branch; only the fidelity and
consistency metrics do, and they do so sharply and consistently across both
neural-network toolkits under identical SMILE settings -- evidence this
reflects the harder, higher-dimensional (30-input) ablation surface these
architectures present to a mask-based surrogate, rather than an
implementation quirk of either framework. On the Geo branch, the two
toolkits' outcomes differ only because of predictive fit, not because of any
difference in GeoSMILE's ability to explain them.

\paragraph{Selection.} On predictive $R^2$, fidelity, consistency, and
stability alone, gradient-boosted trees and Gaussian Process regression are
close: GP wins outright on Feature-branch predictive fit and fidelity, the
boosted trees are more uniform across both branches and cost a fraction of
the compute. Both candidate categories comfortably clear
every guard and show no fidelity collapse, in contrast to the neural
networks (Feature-branch fidelity below $0.13$, and negative once
weighted-adjusted, for both toolkits tested) and to SVR (excluded on Geo
outright). XGBoost was carried forward as the published model on the
strength of the boosting category as a whole -- it is not uniquely superior
to LightGBM or CatBoost on this screen, and Gaussian Process regression
remains a legitimate, quality-competitive alternative that this thesis did
not pursue only because of its computational cost, not because these four
metrics disqualify it. That trade-off, and GP's Feature-branch fidelity
advantage in particular, is noted here as a direction for follow-up work
rather than resolved.
